% Cover letter using letter.sty
\documentclass{letter} % Uses 10pt
%Use \documentstyle[newcent]{letter} for New Century Schoolbook postscript font
% the following commands control the margins:
\topmargin=-0.5in    % Make letterhead start about 1 inch from top of page 
\textheight=8in  % text height can be bigger for a longer letter
\oddsidemargin=0pt % leftmargin is 1 inch
\textwidth=6.5in   % textwidth of 6.5in leaves 1 inch for right margin

\usepackage{graphicx}
\usepackage{wallpaper}
\ULCornerWallPaper{1}{figures/CSM_Letterhead.pdf}

\begin{document}

%\signature{Steven C. DeCaluwe}           % name for signature 
\longindentation=0pt                       % needed to get closing flush left
\let\raggedleft\raggedright                % needed to get date flush left
 
 
\begin{letter}{ }

\vspace{-20mm}
\begin{flushleft}
{\large\bf Steven C. DeCaluwe}
\end{flushleft}
\medskip\hrule height 1pt
\begin{flushright}
\hfill Associate Professor and Director of Graduate Studies\\
\hfill Department of Mechanical Engineering \\
\hfill Colorado School of Mines \\
%\hfill 1601 Illinois St. \\
%\hfill Golden, CO  80401\\
\hfill decaluwe@mines.edu\\
\end{flushright} 
%\vfill % forces letterhead to top of page

\opening{Dear Editors at {\it ACS Energy Letters},} 

\noindent Please find attached our manuscript ``Particle Tracking Methods for Precipitation Reactions in Lithium Batteries,'' submitted for publication in \emph{ACS Energy Letters}. 

The paper explores algorithms for tracking the particle size distribution (PSD) in battery and other energy storage devices where surface-mediated precipitation plays a critical role. Given the importance of such heterogeneous nucleation and growth reactions in next-generation energy storage technologies (including Li-ion batteries, beyond-Li-ion, and critical mineral recovery, for example), accurate modeling tools will be of increasing value in coming years.

In our study, we describe a simple chemical kinetic mechanism for evaluating PSD accuracy and use it to evaluate several common approaches for discretizing and tracking the PSD. We identify a limitation in common discetization, where the PSD is distorted and artificially `flattened' over time due to spatial homogenization within a discretized bin of particle radii. We then propose and demonstrate a coordinate transformation that allows accurate particle size tracking. 

The method is demonstrated for a reduced complexity model, but can readily be applied across a range of arbitrarily complex chemical models.  We anticipate the readership of \emph{ACS Energy Letters} will find the work of great value, and thank you in advance for your consideration.



\closing{Sincerely,\\[-0.5mm]
\fromsig{\includegraphics[scale =0.5, trim = 16mm 6.75in 0 3.4in, clip]{figures/signature.pdf}}\\[1mm]

\fromname{Steven C. DeCaluwe\\
\vspace{3.0mm}
On behalf of:\\
Trent Koberna}
} 





%Steven C. DeCaluwe

%On behalf of:
%Daniel Korff
%Andrew M. Colclasure
%Yeyoung Ha
%Kandler A. Smith
%Trent Koberna



%\noindent 
% \closing{Sincerely yours, 
% \vspace{3mm}
%Enclosures

\end{letter}
 

\end{document}






